% titel: Blatt 7, A3: MO-Diagramm NO, Spin, Konfiguration, Bindungsordnung
% Teil II: der durchdachte Weg.

\subsection*{Gegeben und gesucht}
NO (11 Valenzelektronen). MO-Diagramm, Magnetismus, Elektronenkonfiguration
im Grundzustand, Bindungsordnung von NO und NO$^+$, kurze Bindung.

\subsection*{Ansatz}
11 Elektronen in die MO-Reihenfolge eintragen, ungepaarte z\"ahlen,
Bindungsordnung aus besetzten bindenden und antibindenden Niveaus.

\subsection*{Rechnung}
MO-Energiediagramm (energetisch aufsteigend, mit 1s-Orbitalen als
Rumpf):
\[
  \sigma_{2s}^{2}\ \sigma_{2s}^{*2}\ \pi_{2p}^{4}\ \sigma_{2p}^{2}\
  \pi_{2p}^{*1}
\]
(skizzieren mit Kern-Niveaus 1s ohne Valenzelektronen).

\begin{itemize}
  \item[a)] \textbf{Paramagnetisch:} Ja, es gibt 1 ungepaartes Elektron
        in $\pi^*_{2p}$.
  \item[b)] \textbf{Konfiguration (Grundzustand):}
        \[
          \sigma_{2s}^2\,\sigma_{2s}^{*2}\,\pi_{2p}^4\,\sigma_{2p}^2\,
          \pi_{2p}^{*1}
        \]
  \item[c)] \textbf{Bindungsordnung:}
        \begin{align*}
          \mathrm{BO(NO)} &= \frac{8-3}{2} = 2{,}5,\\
          \mathrm{BO(NO^+)} &= \frac{8-2}{2} = 3{,}0 .
        \end{align*}
        (NO$^+$ $=$ 10 Elektronen, das $\pi^*$-Elektron fehlt:
        $8$ bindende, $2$ antibindende.)
  \item[d)] \textbf{Kurze Bindung:} h\"ohere Bindungsordnung
        $\Rightarrow$ st\"arkere, k\"urzere Bindung. Da
        $\mathrm{BO(NO^+)}=3 > 2{,}5=\mathrm{BO(NO)}$, ist die Bindung in
        \textbf{NO$^+$} k\"urzer.
\end{itemize}

\begin{ergebnis}
NO: paramagnetisch, 1 ungepaartes Elektron; Konfiguration
$\sigma_{2s}^2\sigma_{2s}^{*2}\pi_{2p}^4\sigma_{2p}^2\pi_{2p}^{*1}$.
$\mathrm{BO(NO)}=2{,}5$, $\mathrm{BO(NO^+)}=3{,}0$. Kürzere Bindung in
NO$^+$.
\end{ergebnis}

\begin{kontrolle}
Entfernt man das antibindende Elektron (NO$\rightarrow$NO$^+$), steigt BO
von 2{,}5 auf 3 $\rightarrow$ kürzere/steifere Bindung, chemisch
konsistent. NO ist als Radikal paramagnetisch, NO$^+$ (wie N$_2$+Isolobal)
ist diamagnetisch und stabiler.
\end{kontrolle}

\begin{merksatz}
Ungerade Elektronenzahl $\rightarrow$ ungepaartes MO-Elektron
$\rightarrow$ paramagnetisch. Bindungsordnung = (bindende $-$ antibindende
Elektronen)/2; h\"oher $\Rightarrow$ kürzere Bindung. NO$^+$ ist deshalb
kürzer gebunden als NO.
\end{merksatz}
